\clearpage
\section{Translation}
\label{sec:translation}

The previous sections defined a calculus based on a linear and ordered type
system. In this section we first define an alternative calculus that closer to
how splitting and borrowing of channels would be implemented in a real-world
porgramming language. On its own, the target calculus from this section
provides fewer guarantees. We show that translating a well typed expression
into the target calculus preserves any guarantees made by the much more
expressive type system from before.

\begin{figure}
  \setlength{\abovedisplayskip}{0pt}%
  \setlength{\belowdisplayskip}{0pt}%
  \begin{align*}
    % \TC \grmeq
    %   & \TCUnit
    %     \grmor \TCNew
    %     \grmor \TCFork
    %     \grmor \TCTerm 
    %     \grmor \TCWait
    %     \grmor \TCSend_{\TT}
    %     \grmor \TCRecv_{\TT}
    % \\&
    %     \grmor \TCSplit
    %     \grmor \TCSSplit
    %     \grmor \TCDrop \grmor \TCAcq
    %     \tag{Constants} \\
    \TT \grmeq
        & \TBUnit \grmor \TTArr\TT\TT \grmor \TTProd\TT\TT
                  \grmor \TTSum\TT\TT \grmor \TBChanD \grmor \TBChanF
          \tag{Types}
    \\
    \TCtx \grmeq
        & \TCtxEmpty \grmor \TCtxExt{\TCtx}{\TCtxVar\TEVar\TT}
          \tag{Contexts}
    \\
    \Flag \grmeq
      & \FlagUnset \grmor \FlagSet
    \\
    \TProc \grmeq
        & \TPExp\TE
          \grmor \TPPar\TProc\TProc
          \grmor \TPNew[xx] \TProc
          \grmor \TPSync x \TProc
          \grmor \TPFlag x \Flag
          %\grmor \TPSync* x \TProc
          \tag{Processes}
  \end{align*}
  \caption{Syntax of translated types $\TT$, contexts $\TCtx$, and processes
    $\TProc$}
  \label{fig:unr-syntax}
\end{figure}

\Cref{fig:unr-syntax} defines the syntax of our target calculus. The syntax of
constants, expressions, and values does not changed and is omitted. However, we
view $\TCSend_{\TT}$ and $\TCRecv_{\TT}$ as a family of constants each,
annotated by the message payload type $\TT$.

The type syntax does not differentiate between different kinds of function
arrows. Product types are unordered. $\TBUnit$ and sum types are carried over
unmodified. Channels are untyped and there are now two kinds: data channels
$\TBChanD$ and synchronization channels $\TBChanF$. Communication channels
$\TBChan$ are defined using type abbreviations:
%
\begin{align*}
  \TBChanF* &= \TTSum\TBChanF\TBUnit\\
  \TBChan &= \TTProd{\TBChanF*}{\TTProd{\TBChanD}{\TBChanF*}}
\end{align*}
%
A communication channel is a triple $\targettermcolor{c} =
\TEChan[\TEVar]\TEVary[\TEVarz]$ made up of an optional synchronization channel
$\TEVar$, a data channel $\TEVary$ and a second optional synchronization
channel $\TEVarz$. If $\TEVar$ is a channel, an $\TEApp \TCAcq c$ will have to
be executed first before any communication can happen. Payloads are transferred
via the data channel $\TEVary$. If $\TEVarz$ is an actual channel
$\targettermcolor{c}$ will have to be dropped instead of terminated.

In the syntax of processes, bind groups $\BindGroup$ are eliminated by the
translation. The $\TBindNu$-binder follows the standard syntax more closely.
Additionally, there are $\TBindSync$ and $\TBindSync*$ binders.

\begin{figure}
  \begin{mathpar}
    \RuleUProcRedExp      \and
    \RuleUProcRedFork     \and
    \RuleUProcRedNew      \and
    \RuleUProcRedLSplit   \and
    \RuleUProcRedRSplit   \and
    \RuleUProcRedDrop     \and
    \RuleUProcRedAcquire  \and
    \RuleUProcRedClose    \and
    \RuleUProcRedCom
  \end{mathpar}
  \caption{Reduction of unrestricted processes}
  \label{fig:unr-process-reduction}
\end{figure}

The reduction behaviour of translated expressions is the same as their
untranslated counterpart (\cref{fig:expression-reduction}).
\Cref{fig:unr-process-reduction} defines the reduction relation for translated
processes.

\begin{figure}
  \begin{mathpar}
    \RuleCUUnit \and
    \RuleCUNew  \and
    \RuleCUFork \and
    \RuleCUClose\and
    \RuleCUWait \and
    \RuleCUSend \and
    \RuleCURecv \and
    \RuleCULSplit\and
    \RuleCURSplit\and
    \RuleCUDrop \and
    \RuleCUAcquire
    \\
    \RuleTUConst \and
    \RuleTUVar \and
    \RuleTUAbs
  \end{mathpar}
  etc.---system is standard (no linearity/ordering)
  \begin{mathpar}
    \RuleTProcExpr    \and
    \RuleTProcPar     \and
    \RuleTProcRes     \and
    \RuleTProcSync
  \end{mathpar}
  \caption{Typing of constants, expressions and processes}
  \label{fig:unr-typing}
\end{figure}

\Cref{fig:unr-typing} shows the typing rules for constants, expressions and processes.

\begin{figure}
  \begin{align*}
    \TrT{\BUnit} &= \TBUnit &
    \TrT{\TArr {\Mob} {\Dir,\Eff} {\T[1]} {\T[2]}} &= \TTArr {\TrT{\T[1]}} {\TrT{\T[2]}} \\
    \TrT{\S} &= \TBChan &
    \TrT{\TPair[\Dir] {\T[1]} {\T[2]}} &= \TTProd {\TrT{\T[1]}} {\TrT{\T[2]}} \\
    && \TrT{\TSum {\T[1]} {\T[2]}} &= \TTSum {\TrT{\T[1]}} {\TrT{\T[2]}}
    \\[\smallskipamount]
    \TrCtx{\CtxEmpty} &= \TCtxEmpty &
    \TrCtx{\CtxExt{\Ctx[1]}{\Ctx[2]}} &= \TCtxApp{\TrCtx{\Ctx[1]}}{\TrCtx{\Ctx[2]}} \\
    \TrCtx{\CtxVar\EVar\T} &= \TCtxExt{\TCtxEmpty}{\TCtxVar\EVar{\TrT\T}} &
    \TrCtx{\CtxExtUnord{\Ctx[1]}{\Ctx[2]}} &= \TCtxApp{\TrCtx{\Ctx[1]}}{\TrCtx{\Ctx[2]}}
  \end{align*}
  \caption{Translation of types and contexts}
\end{figure}

\begin{figure}
  \begin{mathpar}
    \RuleTransExpSubVar   \and
    \RuleTransExpCtxVar   \and
    \RuleTransExpAbsLin   \and
    \RuleTransExpSend     \and
    \RuleTransExpRecv     \and
    \RuleTransProcExp   \and
    \RuleTransProcPar   \and
    \RuleTransProcBind  \and
    \RuleTransBindParEmpty  \and
    \RuleTransBindConsSeq   \and
    \RuleTransBindLastSeq
  \end{mathpar}
  \caption{Translation of expressions and processes}
  \label{fig:translation}
\end{figure}

\begin{lemma}[Translation preserves typing]
  \label{lemma:translation-preserves-typing}
%
  For any context $\Ctx$ and substitution $\Sub$ with
  $\TCtx*=\TrCtx\Ctx\setminus\dom\Sub$
%
  \begin{enumerate}
    \item
      If $\HasType[\Eff] \Ctx \E \T$
      and $\HasType \Ctx \E \T \TransExp \TE*$
      then $\HasType {\TCtx*} {\TE*} {\TrT\T}$.
    \item
      If $\ProcHasType \Ctx \Proc$
      and $\ProcHasType \Ctx \Proc \TransProc \TProc*$
      then $\ProcHasType {\TCtx*} {\TProc*}$.
  \end{enumerate}
\end{lemma}

If \cref{lemma:translation-preserves-typing} holds it will also hold for
$\TCtx*=\TrCtx\Ctx$ since the target type system contains rules for weakening.
As written, it guarantees that $\Sub$ has been applied wherever possible.

\begin{lemma}[Bisimulation]
  \label{lemma:translation-bisimulation}
%
  For any context $\Ctx$ and substitution $\Sub$ with
  $\TCtx*=\TrCtx\Ctx\setminus\dom\Sub$
%
  \begin{enumerate}
    \item If
      $\HasType[\Eff[1]] \Ctx {\E[1]} \T$,
      $\HasType \Ctx {\E[1]} \T \TransExp \TE*[1]$,
      and
      $\HasType[\Eff[2]] \Ctx {\E[2]} \T$,
      $\HasType \Ctx {\E[2]} \T \TransExp \TE*[2]$,
      then
      $\E[1] \Reduces \E[2]$
      iff
      $\TE*[1] \Reduces \TE*[2]$.
    \item If
      $\ProcHasType \Ctx \Proc$
      and
      $\ProcHasType \Ctx \Proc \TransProc \TProc*$,
      then for all
      $\ProcHasType \Ctx {\Proc[i]}$
      with
      $\ProcHasType \Ctx {\Proc[i]} \TransProc \TProc*[i]$,
      $\Proc \Reduces \Proc[i]$
      iff
      $\TProc* \Reduces \TProc*[i]$.
  \end{enumerate}
\end{lemma}

\begin{lemma}[Weak subject reduction]
  \label{lemma:translated-subject-reduction}
%
  For any context $\Ctx$ and substitution $\Sub$ with
  $\TCtx*=\TrCtx\Ctx\setminus\dom\Sub$
%
  \begin{enumerate}
    \item If $\HasType[\Eff] \Ctx \E \T$, $\HasType \Ctx \E \T \TransExp \TE*$,
      and $\TE* \Reduces \TE**$, then $\HasType{\TCtx*}{\TE**}{\TrT\T}$
    \item If $\ProcHasType\Ctx\Proc$, $\ProcHasType\Ctx\Proc \TransProc
      \TProc*$, and $\TProc* \Reduces \TProc**$ then
      $\ProcHasType{\TCtx*}{\TProc**}$.
  \end{enumerate}
\end{lemma}
\begin{proof}
  Follows from the bisimulation and typing preserving translation
  (\cref{lemma:translation-preserves-typing}).
\end{proof}
